% Generated from lab2-report.md; SHA256=e02744ce061f4f5b89a6a0538dade9264431f4c823603619dac8602f46e404c3
% -*- coding: utf-8 -*-
% 基于所提供的南开大学实验报告模板；编译器：XeLaTeX。
% 保留提供模板的版心、封面字体、校徽、水印和页眉页脚；正文由Markdown生成。
\documentclass[UTF8,a4paper,10pt]{ctexart}
\usepackage[left=3.17cm,right=3.17cm,top=2.74cm,bottom=2.74cm]{geometry}
\usepackage{amsmath,amssymb,graphicx,float,booktabs,longtable,array,tabularx}
\usepackage{caption,multirow,xcolor,listings,enumitem,needspace}
\usepackage{fancyhdr,tikz,eso-pic}
\usepackage[colorlinks,linkcolor=black,urlcolor=blue,citecolor=black]{hyperref}
% 有 Noto 字体时使用模板字体，否则使用 ctex 默认字体。
\IfFontExistsTF{Noto Serif CJK SC}{\setCJKmainfont{Noto Serif CJK SC}}{}
\newcommand{\yihao}{\fontsize{26pt}{36pt}\selectfont}
\newcommand{\erhao}{\fontsize{22pt}{28pt}\selectfont}
\newcommand{\xiaoer}{\fontsize{18pt}{18pt}\selectfont}
\newcommand{\sanhao}{\fontsize{16pt}{24pt}\selectfont}
\newcommand{\xiaosan}{\fontsize{15pt}{22pt}\selectfont}
\newcommand{\sihao}{\fontsize{14pt}{21pt}\selectfont}
\newcommand{\banxiaosi}{\fontsize{13pt}{19.5pt}\selectfont}
\newcommand{\xiaosi}{\fontsize{12pt}{18pt}\selectfont}
\newcommand{\dawuhao}{\fontsize{11pt}{11pt}\selectfont}
\newcommand{\wuhao}{\fontsize{10.5pt}{15.75pt}\selectfont}
\newcommand{\tabincell}[2]{\begin{tabular}{@{}#1@{}}#2\end{tabular}}
\ctexset{
 section={name={,、},number=\chinese{section}},
 subsection={name={（,）},number=\chinese{subsection}},
 subsubsection={name={,.},number=\arabic{subsubsection}}
}
\pagestyle{fancy}
\fancyhf{}
\fancyhead[L]{\kaishu\small Lab2：物理内存与页表}
\fancyhead[R]{\kaishu\small 操作系统实验报告}
\fancyfoot[C]{\thepage}
\setlength{\headheight}{15pt}
\renewcommand{\headrulewidth}{0.1pt}
\renewcommand{\footrulewidth}{0pt}
\newcommand{\HRule}{\rule{\linewidth}{0.5mm}}
\newcommand{\HRulegrossa}{\rule{\linewidth}{1.2mm}}
\definecolor{nkupurple}{RGB}{112,35,84}
\definecolor{codebg}{RGB}{247,248,250}
\lstset{
 escapeinside={(*@}{@*)},
 breaklines=true,basicstyle=\ttfamily\footnotesize,
 keywordstyle=\color{blue!70}\bfseries,
 commentstyle=\color{gray},stringstyle=\ttfamily,
 columns=fullflexible,keepspaces=true,showstringspaces=false,
 numbers=left,numberstyle=\tiny\color{blue!50},numbersep=6pt,
 frame=trbl,rulecolor=\color{gray!40},backgroundcolor=\color{codebg},
 xleftmargin=1.5em,tabsize=4,aboveskip=0.8em,belowskip=0.8em
}
\setlength{\parindent}{2em}
\setlength{\parskip}{0.25em}
\setlength{\emergencystretch}{3em}
\setlist{nosep,leftmargin=2em}
% 透明紫色水印：12% 不透明度，置于前景以透过代码背景显示。
\newcommand{\watermark}{\AddToShipoutPictureFG{%
 \AtPageLowerLeft{\includegraphics[width=\paperwidth,height=\paperheight]{NKU-watermark.pdf}}}}
\hypersetup{pdftitle={Lab2：物理内存与页表}}

\usepackage{xurl}
\urlstyle{tt}

\IfFontExistsTF{KaiTi}{\setCJKfamilyfont{zhkai}[AutoFakeBold=2]{KaiTi}}{}

\begin{document}
\begin{titlepage}
\centering
\includegraphics[width=0.8\textwidth]{NKU.png}\\[0.8cm]
{\Huge\kaishu\bfseries 南\quad 开\quad 大\quad 学\par}
\vspace{0.9cm}
{\huge\kaishu\bfseries 计\ 算\ 机\ 学\ 院\par}
\vspace{0.7cm}
{\Large\bfseries 操作系统第二次实验报告\par}
\vspace{0.8cm}
\HRule\par\vspace{0.7cm}
{\LARGE\bfseries 物理内存与页表\\[0.3cm]分配算法及 QEMU 调试\par}
\vspace{0.6cm}\HRule\par\vspace{1.2cm}
{\large\kaishu
\begin{tabular}{rl}
组长：&李珉熠\quad 2413982\\[0.5cm]
组员：&吴宇石\quad 2411080\\[0.5cm]
组员：&赵亚鑫\quad 2413715
\end{tabular}\par}
\vfill
{\Large 2026年10月23日\par}
\end{titlepage}
\pagenumbering{Roman}
\renewcommand{\abstractname}{\kaishu\sihao\textbf{摘要}}
\begin{abstract}
我们在最小可执行内核基础上，完成了物理内存探测与页描述符初始化、最佳适应页分配器、伙伴页分配器和简化对象分配器，并通过实际运行检查与课程评分验证了结果。默认物理内存初始化及最佳适应分配评分为20/20；伙伴与对象分配另有独立测试。在分支任务中，我们构建独立带调试信息的QEMU，以客体GDB观察页表和控制寄存器，以宿主GDB跟踪取指及数据读取的地址翻译。实测启动映射采用顶层1GiB叶子，页表查询无需走满三级。我们在报告中说明课程框架、实际修改、AI辅助过程和测试结论。
\par\noindent\textbf{关键词：}物理内存；最佳适应；伙伴系统；对象分配；Sv39；QEMU调试
\end{abstract}
\clearpage
\tableofcontents
\clearpage
\pagenumbering{arabic}
\watermark


\section*{实验基本信息与小组分工}

实验日期：2026-10-23。我们使用的代码来源为：\href{http://oslab.mobisys.top/lab2026/labcode.html}{课程 labcode} 的 lab2.zip。我们保留了课程自带 Best-Fit 两个测试函数和评分测试源码，关键测试结果与调试数值记录在本报告中，完整运行日志保留在本地。


{\small
\setlength{\tabcolsep}{3pt}
\renewcommand{\arraystretch}{1.22}
\begin{longtable}{>{\raggedright\arraybackslash}p{0.25\linewidth}>{\raggedright\arraybackslash}p{0.68\linewidth}}
\toprule
成员 & 模块与报告分工 \\
\midrule\endfirsthead
成员 & 模块与报告分工 \\
\midrule\endhead
2413982 李珉熠（组长） & 环境、物理内存初始化、构建与评分、汇总报告 \\
2411080 吴宇石 & First/Best-Fit、Buddy 与简化 SLUB 算法及测试分析 \\
2413715 赵亚鑫 & 分页切换、内核 GDB 与 QEMU 源码调试、证据整理 \\
\bottomrule
\end{longtable}
}

我们按以上模块组织小组分工，并使用 Codex 作为 AI 辅助工具，协助阅读框架、实现算法、定位错误和整理文档。我们以实际编译、评分和调试输出核对实验结果，并分别理解和复现所负责的内容；报告中的实现与测试结论以实际运行记录为依据。


\section{实验目的}

\begin{enumerate}
\item 从 DTB 获取 RAM 范围，建立 Page 描述符，隔离固件、内核、元数据和 DTB 占用。

\item 理解 First-Fit，完成 Best-Fit；从零实现 Buddy 和简化 SLUB，并验证分配、释放及回收。

\item 理解 Sv39、SATP、叶子 PTE、TLB 与页表遍历；同时调试客体内核与宿主 QEMU。

\end{enumerate}
\section{实验环境与整体逻辑}

我们的实验环境为 Windows 下的 VMware Ubuntu 24.04，RISC-V GCC 15.1.0、GNU make、系统 QEMU 8.2.2、OpenSBI v1.3、gdb-multiarch。我们还在个人目录中构建了 Ubuntu 源码包中的 QEMU 8.2.2 调试版，用于宿主源码调试；构建后直接运行该二进制，没有执行 make install。


\begin{lstlisting}
QEMU 预装固件/镜像 → ROM → OpenSBI → entry.S
→ 建立 Sv39 映射、写 satp、sfence.vma → 高地址 kern_init
→ DTB 解析 → 选择 pmm_manager → page_init → 分配器 check
→ SLUB 对象测试 → 无限循环
\end{lstlisting}

内核链接 VA 起点为 0xffffffffc0200000，加载 PA 为 0x80200000，偏移量为 0xffffffff40000000。VA=PA+offset；PADDR 进行反向转换。在实现和调试中，我们分别区分 Page *（描述符指针）、物理地址 PA、虚拟地址 VA 和物理页号 PPN。


我们的实验目录为 \nolinkurl{/home/lmy/os/lab2}，工具链位于 \nolinkurl{/home/lmy/riscv-elf-toolchains/bin}。


\section{实验内容与实现}

\subsection{物理内存初始化 page\_init}

我们实现了接口 \nolinkurl{static void page_init(void)}。该函数无参数、无返回值，通过全局量 pages、npage 和 va\_pa\_offset 提供初始化结果。


\begin{enumerate}
\item 验证 DTB 的 base、size、加法溢出；当前框架 page2ppn 固定使用 DRAM\_BASE，因此明确拒绝其他基址，避免悄悄错误映射。

\item 设置 va\_pa\_offset，取设备树上界和 PHYSICAL\_MEMORY\_END 的较小值，再向下按 4096 字节对齐。

\item npage 是管理上界的绝对物理页号，不是 RAM 页数；描述符数量为 npage-nbase。

\item 将描述符数组放在链接符号 end 之后并按页向上对齐，先检查整个数组能否放入受管 RAM；每个描述符先标 reserved、ref=0。

\item 空闲起点放在描述符数组之后，向上对齐。保留 OpenSBI、内核及元数据。当前 DTB 在 RAM 顶部附近，为避免覆盖，保守保留从 DTB 所在页至管理上界的尾段；这是安全但可能浪费部分内存的策略。

\item 仅当存在完整空闲页时调用 init\_memmap。框架要求 n>0，因此没有空闲整页时不调用它。

\end{enumerate}
我们在默认 128 MiB RAM 下的实测结果为：管理起点为 \nolinkurl{0x80000000}，上界为 \nolinkurl{0x88000000}，npage-nbase=32768，sizeof(Page)=40；描述符 PA=0x80207000；完整空闲页区间从 \nolinkurl{0x80347000} 到 \nolinkurl{0x87e00000}（不含上界），共 31417 页。64 MiB 实测自动缩小管理上界；256 MiB 实测按框架 0x88000000 上界裁剪。评分脚本关于 npage 的预期针对默认内存，不能把默认 20/20 推广为任意内存配置的评分结论。


\subsection{First-Fit 的理解及改进}

我们首先分析了课程提供的 First-Fit：default\_init 初始化双向链表和空闲计数；init\_memmap 将可用区间组成块，头页 property 保存块长。alloc\_pages 从链表头开始选第一块长度\ensuremath{\geq}n 的块，分割余块；free\_pages 按地址插入，并与前后物理相邻块合并。


我们新增 \nolinkurl{PMM=first} 编译选项，便于现场独立运行课程 First-Fit 管理器。实际启用时修正了参考实现中两处尾插后缺少 break 的问题，保留 basic\_check 和 default\_check 内容不变；First-Fit 及其下层 SLUB 检查已运行通过。


查找时间最坏 O(空闲块数)，地址有序链表便于合并，但不保证最小碎片。可使用分桶、平衡树、位图或最近搜索位置改善搜索成本；也可减少频繁分割。不能保证任何策略都能消除外部碎片。


\subsection{Best-Fit}

我们实现 Best-Fit 时复用框架约定的地址有序空闲链表；遍历所有满足 property\ensuremath{\geq}n 的块，选择 property 最小者。同长时保留首先遇到的低地址块。分配后移除原头、为剩余块设置新头，减少 nr\_free；释放时插入并合并相邻空闲块。init\_pmm\_manager 默认选择 best\_fit\_pmm\_manager。


我们修正了参考链表插入逻辑时发现的尾插入风险：在尾部插入后增加 break，避免循环继续访问刚插入的节点甚至重复插入。除课程原测试外，增加32页隔离池测试：分配32个单页、耗尽返回NULL、释放偶数页形成碎片时不能分配2页、释放奇数页后重新合并成32页，最终计数恢复。


原 basic\_check 与 best\_fit\_check 内容保持不变；新增 best\_fit\_verify 包装调用二者之外的附加测试。普通运行和注入评分都真正使用 Best-Fit。改进可采用按长度索引的树/分桶，同时保留按地址合并结构；当前线性查找为 O(k)，仍可能制造小碎片。


\subsection{Buddy System}

我们新增 buddy\_pmm.c/h，使用0到15阶空闲链表，每阶块大小为2\textasciicircum{}order页。初始化把任意长度初始区间拆成物理页号对齐的二次幂块。申请 n 页向上取整到最小二次幂，从足够大的一阶拆分；伙伴页号由 ppn XOR (1<<order) 得到。释放仅合并同阶、真正的伙伴，不能仅凭地址相邻就合并。


我们的 alloc\_pages(3) 实际占4页；释放时使用原始申请长度3，使 free\_pages 按同样的阶数回收。各阶头标记由 PageProperty 与 property 保存；available 统计实际剩余页数。我们通过 \nolinkurl{make PMM=buddy} 切换管理器，并在切换前执行 make clean，避免复用另一配置的目标文件。


测试覆盖：零长度和超最大阶请求、3页向4页取整、4页/8页对齐、并存块无重叠、释放前后计数恢复；隔离一个8页块，拆成8个单页至耗尽，逆序释放，重新申请8页必须恢复原块。这能检查分裂、伙伴合并和计数，而不仅是“申请返回非空”。


我们还增加了16页隔离测试：先申请物理对齐的16页块并暂存其余空闲链表，清理测试块管理状态，将第7号页设为保留，只把其余15页加入测试池。逐页申请15次，逐项检查范围、不含保留页、无重复；第16次申请应失败。释放15页后解除中间页保留并归还，检查合并为4阶16页块。再申请3页和5页，分别占4页和8页，剩余4页；释放后恢复16页，重新申请整块成功，最后恢复原链表与全局空闲计数。我们通过手动GDB断点核对到空闲数15、0、16、4、16，以及全局计数恢复。


\subsection{简化 SLUB 对象分配}

我们新增 kmalloc(size\_t n)、kfree(void *p)、slub\_check()。这是体现 SLUB 主体思想的单CPU教学实现，不等同于 Linux 生产级 SLUB。


下层使用当前 pmm\_manager 的页分配；上层提供16、32、64、128、256、512、1024字节七种缓存。每个slab一页，页首为元数据，后面切分为固定长度对象；空闲对象首部存放下一个空闲对象指针。对象16字节对齐。超过1024字节走连续页大对象路径，申请页数为 ceil((header+n)/4096)，最大接受1 MiB。大对象仍保留首部元数据，使 kfree 能还原页及长度。


kfree(NULL) 无操作。释放小对象时检查页头、位置、对齐及重复释放，把对象放回空闲链；整slab空闲后从缓存链移除并归还底层页。Buddy 模式下底层页请求可能取整，但仍保存原请求页数，以相同请求释放。


测试覆盖各大小类边界1/16/17/32/33/.../1024/1025、4096和9000字节，逐字节写入和核对防止重叠，逆序释放后空闲页计数恢复；300个32字节对象跨多个slab，验证数据、同一可用slab中的对象复用及所有页回收。我们在早期测试中错误假定任意slab释放的对象都是下次首选；核对缓存链选择顺序后，我们改为释放当前优先缓存slab中的对象再检验复用。


\subsection{任务规格与 AI 辅助过程}

我们在持续协作中使用 Codex 辅助阅读、编码、排错和解释，围绕接口、约束与验收标准组织任务。以下提示词是我们对协作需求的整理，便于复现实现和测试思路。


\textbf{物理页初始化与 Best-Fit。} 我们要求基于 DTB 初始化 pages、npage 和 va\_pa\_offset，检查溢出及内存边界，保留固件、内核、描述符和 DTB 占用；Best-Fit 遍历全部合适块选择最小者，保持地址有序及相邻块合并。我们保留课程 basic\_check、best\_fit\_check 与评分测试，另加耗尽、碎片和合并测试。


\textbf{Buddy 实现提示。} “实现独立 pmm\_manager：0到15阶，每块为2\textasciicircum{}order页；初始化按绝对物理页号对齐分解区间；申请向上取整并逐阶拆分，释放按 ppn XOR (1<<order) 寻找同阶伙伴并递归合并。支持 PMM=buddy，明确原始申请长度的释放约定。”


\textbf{Buddy 测试提示。} “覆盖零与过大请求、3页申请实际占4页、对齐和不重叠；隔离8页，拆成8个单页至耗尽，逆序释放后重新申请完整8页，恢复原状态并核对空闲计数。”


\textbf{SLUB 实现提示。} “在页分配器上实现 kmalloc/kfree，提供16到1024字节七种缓存，每slab一页、对象16字节对齐、空闲对象内嵌next指针；大对象以带头部的连续页分配，限制1MiB，完全空闲时归还页，兼容Best-Fit和Buddy。”


\textbf{SLUB 测试提示。} “覆盖大小类边界、零与过大请求、跨页大对象；逐字节检查数据，300个32字节对象跨多个slab，检查当前可用slab内的对象复用，全部释放后核对底层页数，在两种页分配器下分别运行。”


\textbf{QEMU 调试提示。} “在个人目录构建带符号的QEMU，用宿主GDB运行模拟器，客体GDB连接stub；分别观察kern\_init取指和boot\_dtb数据读，检查TLB miss、页表项读取与最终物理地址，排除is\_debug访问，并对照Bare与Sv39。”


我们通过实际构建与运行核对 AI 辅助生成的代码，发现不符合框架或测试前提的内容后继续修改；测试结果以第五节的实际验证范围为准。


我们在实验中实际修复的问题包括：QEMU 通用loader缺少入口交接信息，改用-kernel；入口缺失低地址过渡映射，增加顶层第2项恒等映射；工具链明确rv64gc/lp64d并改正旧sfence.vm别名；课程脚本CRLF在sh下报错，统一LF；Best-Fit尾插入增加break；纠正SLUB复用测试前提；QEMU源码包缺少可安装BIOS资源时使用disable-install-blobs，运行时-L /usr/share/qemu复用系统固件。


\section{分支任务：实际调试 QEMU 的 MMU 路径}

\subsection{调试架构}

我们用客体 gdb-multiarch 经 QEMU stub 观察 RISC-V 寄存器和内核，用宿主普通 gdb 运行 x86-64 的 QEMU 进程，并在其C源码中设置断点。宿主GDB作为QEMU父进程，不需要sudo attach，也避免ptrace限制。QEMU调试版位于 /home/lmy/oslab-tools/qemu-ubuntu-build/qemu-system-riscv64，源码在相邻qemu-source。我们分别检查了内核启动、页表设置、取指和数据读取的翻译路径。


\subsection{SATP 与大页的客体观察}

我们在客体 GDB 中读到 SATP=0x8000000000080205：MODE=8表示Sv39，ASID=0，根表PPN=0x80205，即物理地址0x80205000。根表第511项位于0x80205ff8，PTE=0x200000cf；恒等映射第2项也为相同PTE。


PTE的PPN=0x80000，对应PA基址0x80000000；低位0xcf表示V/R/W/X/A/D为1、U/G为0。根表叶子映射1 GiB。三级页表是最多三级查询，遇到叶子可提前结束。本实验两项覆盖低地址0x80000000起1 GiB，以及高地址0xffffffffc0000000起1 GiB。真实受管RAM只有其中一部分，大页映射范围不等于可分配RAM范围。


\subsection{宿主断点与真实翻译}

我们分别选择 kern\_init 的取指和 boot\_dtb 的数据读取做两次独立运行。宿主条件断点位于 riscv\_cpu\_tlb\_fill、get\_physical\_address、PTE读取后的分支和tlb\_set\_page前。is\_debug=0用于排除调试器自身的内存查询。


我们记录到以下关键输出：


\begin{lstlisting}
WALK_ENTER addr=0x80200000 satp=0 is_debug=0
WALK_ENTER addr=0xffffffffc02000d6 satp=0x8000000000080205 is_debug=0
PTE_READ ... i=0 ptshift=18 base=0x80205000 pte_addr=0x80205ff8 pte=0x200000cf
TRANSLATED VA=0xffffffffc02000d6 PA=0x802000d6 prot=7 cached_size=0x1000
TLB_MISS address=0xffffffffc0206008 access=0 mmu_idx=1
TRANSLATED VA=0xffffffffc0206008 PA=0x80206008 prot=7 cached_size=0x1000
\end{lstlisting}

取指access=2，数据读取access=0，mmu\_idx=1对应此次S模式访问。数据读调用栈真实出现 do\_ld8\_mmu \ensuremath{\rightarrow} mmu\_lookup \ensuremath{\rightarrow} mmu\_lookup1 \ensuremath{\rightarrow} tlb\_fill \ensuremath{\rightarrow} riscv\_cpu\_tlb\_fill。


\subsection{页表遍历与TLB原理}

QEMU 8.2的 target/riscv/cpu\_helper.c 中 get\_physical\_address 先检查权限模式和SATP模式；本次Bare阶段直接使用原地址。Sv39时以根表物理地址为base，计算当前层index和pte\_addr=base+index*8，使用address\_space\_ldq读取8字节PTE，再检查合法性。V=0失败；R/W/X有叶子权限则转到leaf；否则提取PPN，将base更新为下级表物理地址，继续循环。因此循环的职责是层级遍历，不是重复猜地址。框架顶层即为叶子，所以这次循环i=0时结束。


此处ptshift=18是还要保留的下级VPN位数，两层各9位，共18位；不是虚拟地址总偏移。1 GiB叶子的物理地址由PPN高位、虚拟地址的低级VPN和12位字节偏移组成，等价于 PA=0x80000000+(VA \& 0x3fffffff)。对数据读得到0x80206008。


软件TLB查找位于accel/tcg/cputlb.c。mmu\_lookup1读取并比较TLB标签，不命中时先尝试victim\_tlb\_hit，再通过tlb\_fill调用架构的riscv\_cpu\_tlb\_fill；成功后tlb\_set\_page缓存结果。取指另由get\_page\_addr\_code\_hostp/probe\_access\_internal路径进入。get\_physical\_address不是每次访问必经函数：命中缓存时可以绕过页表遍历，TCG翻译块缓存也会影响取指路径。


QEMU软件TLB承担客体地址转换及访问权限缓存，并保存加速宿主访问需要的信息；它不等于实体RISC-V处理器某一型号的硬件TLB，也不等于宿主x86 CPU自己的TLB。本次大页翻译虽然使用1 GiB叶子，软件缓存填充实测cached\_size为0x1000，不能把二者混为相同页大小。进入内核前Bare模式虽然不做Sv39遍历，QEMU仍可能调用软件TLB填充，这正是日志中SATP=0时也有TLB\_MISS的原因。


\section{测试结果}

{\small
\setlength{\tabcolsep}{3pt}
\renewcommand{\arraystretch}{1.22}
\begin{longtable}{>{\raggedright\arraybackslash}p{0.25\linewidth}>{\raggedright\arraybackslash}p{0.29\linewidth}>{\raggedright\arraybackslash}p{0.37\linewidth}}
\toprule
测试 & 实测结果 & 验证内容 \\
\midrule\endfirsthead
测试 & 实测结果 & 验证内容 \\
\midrule\endhead
课程 page\_init 与 Best-Fit 评分 & 20/20 & 默认128MiB配置下的课程评分 \\
Best-Fit原测试与附加32页碎片/合并 & PASS & 耗尽、连续性、释放后合并和计数恢复 \\
First-Fit课程自检与SLUB & PASS & 独立构建运行、页分配与对象回收 \\
Buddy取整/对齐/耗尽/合并 & PASS & 8页拆分；16页保留页、唯一性、3/5页取整与全量恢复 \\
SLUB尺寸边界/数据/多slab/复用/回收 & Best-Fit、Buddy两种下层均PASS & 数据不重叠、对象复用、底层页归还 \\
64MiB、256MiB启动与内存边界 & 启动及测试成功 & 64MiB自动缩小，256MiB按128MiB上界裁剪 \\
客体GDB与SATP/PTE & 高地址断点及根表值已确认 & SATP=0x8000000000080205，PTE=0x200000cf \\
宿主QEMU取指与数据读翻译 & 条件断点命中并确认PA & 分别得到PA=0x802000d6与0x80206008 \\
\bottomrule
\end{longtable}
}

我们在限时验证中使用 timeout 结束运行，在宿主跟踪完成后由调试脚本主动结束 QEMU，因此相应的 timeout 和 killed 信息属于验证流程结束。我们在本地保留了真实文本日志；测试范围为上表所列场景，尚未覆盖生产级并发、任意恶意 DTB 和全部硬件配置。


\section{总结、理论联系与扩展思考}

\subsection{内存管理各层的关系}

页级分配器解决“哪些连续物理页可以使用”，SLUB解决“如何把页切成对象”，页表解决“CPU地址如何翻译”，TLB解决“如何缓存翻译”；这些是不同层次。reserved防止把固件/内核/描述符作为空闲内存发出；property只在空闲块头解释，不能拿任意页的值当块长。


\subsection{扩展任务 Challenge：硬件的可用物理内存范围的获取方法（思考题）}

\textbf{问题：如果 OS 无法提前知道当前硬件的可用物理内存范围，有何办法让 OS 获取可用物理内存范围？}


我们的思路是由内核在启动时读取固件或启动程序提供的硬件描述，而不是把某一种机器的内存容量写死。可用的信息来源包括设备树 DTB、UEFI 内存映射或平台专用固件协议。DTB 是设备树的二进制数据，用于描述硬件；UEFI 是一种固件接口，其内存映射可描述各内存区域的类型和用途。


具体处理分为三步：


\begin{enumerate}
\item 获取内存布局：启动程序传入硬件描述的位置，内核解析其中的 RAM 起始地址、大小以及可获得的区域属性，检查边界和地址计算是否溢出。

\item 确定可分配范围：从 RAM 中排除固件、内核、启动栈、页表、页描述符数组、设备树以及其他保留区域；设备的地址窗口也不能作为普通空闲 RAM 分配。因此，硬件提供的 RAM 总范围不等于当前可分配的空闲范围。

\item 建立管理结构：将剩余区间的起点向上、终点向下按页对齐，为完整物理页建立管理状态，再交给页分配器。若硬件存在多段不连续 RAM，应分别记录，不能把中间的地址空洞当作内存。

\end{enumerate}
在我们本次实验中，启动环境通过 a1 寄存器传入 DTB 的物理地址，内核解析 memory 节点的 reg 属性，取得 RAM 起点 0x80000000 和大小 128 MiB。我们再结合内核布局计算空闲范围，最终将 [0x80347000,0x87e00000) 的完整页交给分配器。


如果没有标准硬件描述，我们可以考虑使用平台专用的内存控制器查询接口，或在已知平台规则和异常处理机制下进行受控探测。但不能盲目向任意物理地址读写测试，否则可能覆盖已有数据、访问设备寄存器或触发访问异常。


我们当前的 DTB 解析只识别一个 memory 节点中的一段内存范围，尚未完整处理多段 RAM、设备树保留表和 reserved-memory 节点。后续可扩展这些信息的解析，建立更完整的可用与保留区间列表。


\subsection{当前实现的局限与总结}

我们本次尚未实现进程地址空间隔离、按需分页、页面置换、系统调用、调度、文件系统，以及多核/中断环境下的分配器同步。过渡恒等映射和宽权限1GiB启动页适合教学启动，不是最终内核应长期采用的最小权限布局。


通过本次实验，我们完成了可管理物理页的建立、Best-Fit接入、Buddy和对象分配，并在QEMU源码中观察了真实地址转换的关键路径。


参考：\href{http://8.135.34.58/lab2026/_book/lab2/lab2_2_1_exercise.html}{课程练习}、\href{http://8.135.34.58/lab2026/_book/lab2/lab2_4_qemu_mmu_debug.html}{分支任务}、\href{https://github.com/qemu/qemu/blob/v8.2.2/target/riscv/cpu_helper.c}{QEMU v8.2.2 MMU源码}、\href{https://github.com/qemu/qemu/blob/v8.2.2/accel/tcg/cputlb.c}{QEMU软件TLB}。我们实际调试时采用Ubuntu重打包的8.2.2源码。


\section{关键代码的调试验证}

我们进一步通过实际\nolinkurl{make gdb}入口检查启动状态，并用断点观察分配算法的区别。\nolinkurl{make gdb}内部调用工具链GDB加载符号、设置RISC-V架构并连接1234端口；我们读到的SATP和PTE与前述数值一致。


我们在best\_fit\_check第237行申请1页之前，通过GDB确认两个空闲块相对p0偏移分别为1和4，长度分别为2和1，总空闲3页；执行该行后p1-p0=4、剩余2页；接着第238行申请连续2页仍成功，空闲页数归零。因此第239行assert(p0 + 4 == p1)结合前一行连续申请，明确验证最小合适块选择，与First-Fit的首次合适块选择不同。我们还通过GDB确认描述符大小为40字节，page\_link偏移为24字节。


Page管理标记不是PTE硬件权限；单张Sv39表为4KiB、512项，而完整多级结构按实际映射分配；启动根表使用1GiB叶子所以不走满三级。测试断言属于覆盖特定行为的证据，不能代替任意非法释放或并发正确性的证明。这些观察帮助我们把分配策略、链表结构和实际执行结果对应起来。


\end{document}
